\documentclass[10pt,a4paper]{amsart}
\usepackage[margin=19mm]{geometry}
\usepackage{amsmath,amssymb,mathtools}
\usepackage[scr=rsfs,cal=euler]{mathalfa}
\usepackage{xfrac}
\usepackage[unicode,hidelinks,bookmarks=false]{hyperref}
\newcommand{\ie}{i.e.\ }
\newcommand{\cf}{cf.\ }
\newcommand{\resp}{respectively\ }
\newcommand{\Subsection}{\subsection}
\providecommand{\nobreakdash}{\nobreak\hbox{-}}
\setlength{\emergencystretch}{2em}
\multlinegap0pt
\usepackage{amssymb}
\usepackage{mathtools}
\usepackage{enumerate}
\usepackage{dsfont}
\usepackage{tikz}
\newtheorem{proposition}{Proposition}
\def\N{\mathbb{N}}
\def\Q{\mathbb{Q}}
\newcommand{\Set}[1]{\left\{ #1 \right\}}
\def\Z{\mathbb{Z}}
\newcommand{\brac}[1]{\left({#1}\right)}
\def\lcm{\operatorname{lcm}}
\def\mJ{\mathcal{J}}
\newcommand{\set}[1]{\{ #1 \}}
\title{Smoothness of Markov Partitions for Expanding Toral Endomorphisms}
\author{Chayce Hughes, Huub de Jong}
\date{Source: arXiv:2604.00257v1 (CC BY 4.0)}
\begin{document}
\maketitle

\noindent\emph{Source and licence.} Smoothness of Markov Partitions for Expanding Toral Endomorphisms. Version arXiv:2604.00257v1 (CC BY 4.0), distributed by arXiv under the Creative Commons Attribution 4.0 International licence, http://creativecommons.org/licenses/by/4.0/, as recorded in the arXiv licence metadata for this exact version and shown on the abstract page https://arxiv.org/abs/2604.00257v1. Full licence notice: \texttt{licenses/arxiv-2604.00257-CC-BY-4.0.txt}.\par\smallskip

\noindent\textit{Editorial scope.} Counted unit: the article's proposition prop:bounded\_power (source0.tex lines 353--355). The article's complete original proof of it is delivered with it (source0.tex lines 357--384, one \texttt{proof} environment of 28 lines). No mathematics is written, repaired or re-derived: the statement and the proof are the article's own lines, in the article's own notation, and no retained line is substituted or re-flowed.  No further source-local context is needed: every cross-reference of every retained line resolves inside the two delivered spans.  Nothing was omitted from any retained span, so the package carries no omission record.  No retained line contains a citation, so the wrapper carries no bibliography and no bibliography entry.  No retained line contains a figure environment, an image-inclusion command or drawing code, so no image is delivered and the package declares no asset dependency.  Editorial formatting only, in addition: an AMS \texttt{amsart} preamble that restates the article's own declarations and macros used by the retained lines, namely the environments \texttt{proposition} on separate counters, together with the standard packages those lines need.  The retained environments print in this document as Proposition 1 in order of appearance, whereas the article prints the same items with the numbering of its own full document.  The complete native source of the article as distributed in the arXiv e-print for this exact version is bundled unchanged as \texttt{sources/197580/source0.tex}.
\begin{proposition}\label{prop:bounded_power}
    For all $d\in \N$, $\mathcal{L}(d)$ is finite.
\end{proposition}

\begin{proof}
    Let $A\in \mathcal{A}_d$ and $\lambda_j = r_je^{i\theta_j}$ with $r_j \ge 0$ denote its eigenvalues. Observe that if the characteristic polynomial $\chi_A(x)$ is a product of two polynomials $f(x), g(x) \in \Q[x]$, then $f$ and $g$ partition the roots of $\chi_A$. Hence, $n_A$ will be the least common multiple of the smallest $n_f$ and smallest $n_g$ for which $\xi^{n_f} \in \Z$ and $\omega^{n_g} \in \Z$ for all roots $\xi$ of $f$ and $\omega$ of $g$. Since $\chi_A$ consists of at most $d$ irreducible factors, it thus suffices to bound $n_A$ by a function of $d$ when $\chi_A$ is irreducible over $\Q$.
    
    Assuming $\chi_A$ is irreducible, write $k = \lambda^{n_A} \in \Z$ for some root $\lambda$ of $\chi_A$. Then, $\lambda$ is a root of $x^{n_A} - k$, so $\chi_A$ divides $x^{n_A} - k$. This means all the roots of $\chi_A$ have the same modulus, so $r_1=r_2=...=r_d=r$ for some $r\geq0$. When $r=0$ the situation is of course trivial, so we consider $r>0$. 
    
    
    Furthermore, the constant coefficient of $\chi_A$ is $\prod_{j=1}^d \lambda_j\in \Z$, so $r^d \in \N$. Thus, the smallest $\nu \in \N$ for which $r^\nu \in \N$ is at most $d$. After raising $A$ to the power $\nu$, and taking its characteristic polynomial, we obtain 
    $$
        \chi_{A^\nu}( r^\nu x) = \prod_{j=1}^d \brac{r^\nu x - r^\nu e^{i \nu\theta_j}}= r^{\nu d} \underbrace{\prod_{j=1}^d \brac{x - e^{i \nu\theta_j}}}_{\in \Q[x]}.
    $$
    Let $\varphi$ denote Euler's totient function and $\Phi_m$ the $m^{\text{th}}$ cyclotomic polynomial, so that $\deg(\Phi_m) = \varphi(m)$. By assumption, $e^{in_A\theta_j} = \pm 1$, so $\theta_j \in \pi \Q$. That is, all the roots of $\chi_{A^\nu}( r^\nu x)$ are roots of unity. Hence, $\chi_{A^\nu}( r^\nu x)/r^{\nu d}$ is a product of cyclotomic polynomials:
    $$
        \chi_{A^\nu}( r^\nu x) = r^{\nu d} \prod_{j=1}^d \brac{x - e^{i \nu\theta_j}} = r^{\nu d} \prod_{J \in \mJ} \Phi_{m_J}(x)
    $$
    for some indexing set $\set{m_J: J \in \mJ} \subset \N$. The smallest natural $n_J$ such that $\xi^{n_J} =\pm 1$ for all roots $\xi$ of $\Phi_{m_J}$ is
    $$
        n_J = \min \Set{t \in \N: \frac{2\pi t}{m_J} = 0 \mod \pi} = \frac{m_J}{\gcd(2,m_J)} \le m_J.
    $$
    Additionally, $\deg(\Phi_{m_J}) = \varphi(m_J) \le d$ since $\deg(\chi_{A^\nu}) = d$. It is well known that $\lim_{t \to \infty} \varphi(t) = \infty$, so there are only finitely many cyclotomic polynomials with degree at most $d$:
    $$
         m_J \le \max_{1 \le t \le d} \max \varphi^{-1}(t) < \infty
    $$
    Thus, we observe that
    $$
        n_A = \nu \, \lcm\set{n_J : J \in \mJ} \le d \prod_{J} n_J \le d \prod_{J} m_J \le d \brac{\max_{1 \le t \le d} \max \varphi^{-1}(t)}^d < \infty,
    $$
    as desired.
\end{proof}

\end{document}
